\chapter{System Design}

\section{System Architecture}
The system adopts a decoupled Client-Server architecture. The Presentation Layer is built using Streamlit, communicating with a Flask RESTful API Backend via HTTP protocols. Data is stored in an SQLite relational database alongside a localized media store.

\begin{figure}[H]
\centering
\begin{tikzpicture}[
    node distance=1.5cm and 2cm,
    block/.style={rectangle, draw=blue!70, fill=blue!5, rounded corners, minimum width=3cm, minimum height=1cm, align=center, font=\sffamily\small},
    backend/.style={rectangle, draw=green!70, fill=green!5, rounded corners, minimum width=3cm, minimum height=1cm, align=center, font=\sffamily\small},
    db/.style={cylinder, draw=orange!70, fill=orange!5, shape border rotate=90, aspect=0.25, minimum width=2.2cm, minimum height=1.2cm, align=center, font=\sffamily\small},
    line/.style={draw, -Latex, thick, font=\sffamily\scriptsize}
]

% Frontend Nodes
\node[block] (user) {User Portal\\(\texttt{user\_front.py})};
\node[block, below=0.5cm of user] (judge) {Judge Portal\\(\texttt{judge\_front.py})};
\node[block, below=0.5cm of judge] (admin) {Admin Portal\\(\texttt{admin\_front.py})};

% Subsystem Box Frontend
\begin{scope}[on background layer]
    \node[draw=blue!40, fill=blue!2, dashed, fit=(user) (judge) (admin), inner sep=0.4cm, label=above:\textbf{Streamlit Frontend Layer}] (fe_box) {};
\end{scope}

% Backend Nodes
\node[backend, right=3.5cm of judge] (api) {Flask REST API\\(\texttt{app.py}:9999)};
\node[backend, above=0.8cm of api] (ai) {AI Validation Module\\(Image Classifier)};

% Storage Nodes
\node[db, right=3cm of api] (database) {SQLite DB\\(\texttt{submissions.db})};
\node[block, draw=red!70, fill=red!5, below=0.8cm of database] (storage) {Local File Store\\(\texttt{/uploads})};

% Connections
\path [line] (user.east) -- node[above, sloped] {POST /submissions} (api.west);
\path [line] (judge.east) -- node[above, sloped] {PUT /submissions/\{id\}} (api.west);
\path [line] (admin.east) -- node[below, sloped] {GET, DELETE} (api.west);

\path [line] (api.north) -- node[right] {Validate Image} (ai.south);
\path [line] (api.east) -- node[above] {SQL Queries} (database.west);
\path [line] (api.east) -- node[below, sloped] {Save File} (storage.west);

\end{tikzpicture}
\caption{System Architecture Diagram}
\label{fig:sys_arch}
\end{figure}

\section{Database Design (Entity-Relationship Diagram)}
The relational schema comprises three main entities: \texttt{USERS}, \texttt{SUBMISSIONS}, and \texttt{EVALUATIONS}.

\begin{figure}[H]
\centering
\begin{tikzpicture}[
    entity/.style={rectangle, draw=black, fill=gray!10, thick, minimum width=2.8cm, minimum height=1.2cm, align=center, font=\sffamily\bfseries},
    relationship/.style={diamond, draw=black, fill=yellow!10, thick, aspect=1.5, minimum width=2cm, align=center, font=\sffamily\small},
    line/.style={draw, thick}
]

% Entities
\node[entity] (user) {USERS\\{\normalfont\tiny (id, name, email, role)}};
\node[relationship, right=2.5cm of user] (submits) {Submits};
\node[entity, right=2.5cm of submits] (submission) {SUBMISSIONS\\{\normalfont\tiny (id, title, image\_path, film\_stock, camera, status)}};
\node[relationship, below=2cm of submission] (evaluates) {Evaluates};
\node[entity, left=2.5cm of evaluates] (judge) {JUDGES\\{\normalfont\tiny (id, name, expertise)}};

% Connections
\path [line] (user) -- node[above] {1} (submits);
\path [line] (submits) -- node[above] {N} (submission);
\path [line] (judge) -- node[above] {1} (evaluates);
\path [line] (evaluates) -- node[right] {N} (submission);

\end{tikzpicture}
\caption{Entity-Relationship Diagram (ERD)}
\label{fig:erd}
\end{figure}

\section{Sequence Diagram: Submission Workflow}
The sequence below illustrates client-side image compression followed by asynchronous REST interaction with the Flask Backend.

\begin{figure}[H]
\centering
\begin{tikzpicture}[
    actor/.style={rectangle, draw, fill=blue!10, minimum width=2cm, minimum height=0.8cm, align=center, font=\sffamily\small},
    lifeline/.style={draw, dashed, gray},
    msg/.style={draw, -Latex, thick, font=\sffamily\scriptsize},
    reply/.style={draw, -Latex, dashed, thick, font=\sffamily\scriptsize}
]

% Lifeline Headers
\node[actor] (c) {Contestant};
\node[actor, right=2cm of c] (fe) {Frontend\\(Streamlit)};
\node[actor, right=2.5cm of fe] (be) {Backend\\(Flask API)};
\node[actor, right=2.5cm of be] (db) {Database\\(SQLite)};

% Vertical Lifelines
\draw[lifeline] (c) -- ++(0,-6.5);
\draw[lifeline] (fe) -- ++(0,-6.5);
\draw[lifeline] (be) -- ++(0,-6.5);
\draw[lifeline] (db) -- ++(0,-6.5);

% Messages
\draw[msg] ([yshift=-1cm]c.south) -- node[above] {1. Select image \& fill form} ([yshift=-1cm]fe.south);
\draw[msg] ([yshift=-1.8cm]fe.south) -- node[above] {2. Compress Image (PIL)} ([yshift=-1.8cm]fe.south);
\draw[msg] ([yshift=-2.6cm]fe.south) -- node[above] {3. POST /api/submissions (FormData)} ([yshift=-2.6cm]be.south);

\draw[msg] ([yshift=-3.4cm]be.south) -- node[above] {4. Save file to /uploads} ([yshift=-3.4cm]be.south);
\draw[msg] ([yshift=-4.2cm]be.south) -- node[above] {5. INSERT INTO submissions} ([yshift=-4.2cm]db.south);
\draw[reply] ([yshift=-5.0cm]db.south) -- node[above] {6. Return DB Record ID} ([yshift=-5.0cm]be.south);

\draw[reply] ([yshift=-5.8cm]be.south) -- node[above] {7. HTTP 201 Created (JSON)} ([yshift=-5.8cm]fe.south);
\draw[reply] ([yshift=-6.3cm]fe.south) -- node[above] {8. Render Success Balloon UI} ([yshift=-6.3cm]c.south);

\end{tikzpicture}
\caption{Sequence Diagram for Image Submission}
\label{fig:seq_diagram}
\end{figure}
\section{Database Design (Entity-Relationship Diagram)}
The SQLite database consists of three main entities: \texttt{USERS}, \texttt{SUBMISSIONS}, and \texttt{EVALUATIONS}. The diagram below illustrates the relational table structures, primary keys (PK), foreign keys (FK), and their cardinality relationships.

\begin{figure}[H]
\centering
\begin{tikzpicture}[
    node distance=2.5cm and 2cm,
    table/.style={
        rectangle split,
        rectangle split parts=2,
        draw=blue!80!black,
        fill=blue!5,
        rounded corners=2pt,
        minimum width=3.5cm,
        font=\sffamily\small,
        align=left
    },
    tableheader/.style={fill=blue!20, font=\sffamily\bfseries\small},
    line/.style={draw=black!70, -Latex, thick, font=\sffamily\scriptsize}
]

% USERS Table
\node[table] (users) {
    \textbf{USERS}
    \nodepart{second}
    \underline{id} (PK) : INTEGER \\
    name : VARCHAR(100) \\
    email : VARCHAR(100) \\
    role : VARCHAR(20)
};

% SUBMISSIONS Table
\node[table, right=2.5cm of users] (submissions) {
    \textbf{SUBMISSIONS}
    \nodepart{second}
    \underline{id} (PK) : INTEGER \\
    user\_id (FK) : INTEGER \\
    title : VARCHAR(150) \\
    film\_stock : VARCHAR(50) \\
    camera : VARCHAR(50) \\
    format : VARCHAR(20) \\
    image\_path : VARCHAR(255) \\
    status : VARCHAR(20) \\
    created\_at : DATETIME
};

% EVALUATIONS Table
\node[table, below=2cm of submissions] (evaluations) {
    \textbf{EVALUATIONS}
    \nodepart{second}
    \underline{id} (PK) : INTEGER \\
    submission\_id (FK) : INTEGER \\
    judge\_id (FK) : INTEGER \\
    score : FLOAT (1--10) \\
    feedback : TEXT \\
    created\_at : DATETIME
};

% Relationships / Connectors
\path[line] (users.east) -- node[above] {1 : N (Submits)} (submissions.west);
\path[line] (submissions.south) -- node[right] {1 : N (Has)} (evaluations.north);
\path[line] (users.south) |- node[below, near start] {1 : N (Evaluates as Judge)} (evaluations.west);

\end{tikzpicture}
\caption{Entity-Relationship Diagram (Database Relational Schema)}
\label{fig:erd_detailed}
\end{figure}